% ================================================================ chapter 4
\chapter{Experimental Setup}
\label{ch:setup}

\section{Datasets}
Table~\ref{tab:datasets} lists the datasets used. FakeAVCeleb%
~\cite{khalid2021fakeavceleb} provides native four-quadrant labels and is the
primary training set. AV-Deepfake1M~\cite{cai2023avdeepfake1m} and
LAV-DF~\cite{cai2022lavdf} drive the temporal-localization task. DFDC%
~\cite{dolhansky2020dfdc} and KoDF are reserved exclusively for cross-dataset
evaluation. ASVspoof~\cite{wang2020asvspoof} strengthens the audio branch.

\begin{table}[h]
\centering
\caption{Datasets used in this study.}
\label{tab:datasets}
\small
\begin{tabular}{llll}
\toprule
\textbf{Dataset} & \textbf{Modality} & \textbf{Labels} & \textbf{Role} \\
\midrule
FakeAVCeleb   & A+V & 4-quadrant           & train + in-domain test \\
AV-Deepfake1M & A+V & clip + localization   & train + localization \\
LAV-DF        & A+V & localized forgeries   & localization \\
DFDC          & A+V & binary                & cross-dataset test only \\
KoDF          & A+V & binary                & cross-dataset test only \\
ASVspoof 2019/21 & A & bona fide / spoof    & audio-branch auxiliary \\
\bottomrule
\end{tabular}
\end{table}

\section{Preprocessing}
Video is decoded at 25\,fps; faces are detected, tracked, and aligned, producing a
$224\times224$ face crop and a $96\times96$ mouth crop. Audio is resampled to
16\,kHz mono and loudness-normalized. All datasets are converted to a unified
manifest schema with per-modality labels, quadrant, manipulated intervals,
generator identity, and demographic metadata. Splits are subject-disjoint;
leave-one-generator-out (LOGO) splits hold out one manipulation method entirely.

\section{Implementation Details}
The framework is implemented in PyTorch. Experiments run on an NVIDIA DGX Spark
system (GB10 Grace--Blackwell, 128\,GB unified memory) using BF16 precision.
Training follows a cached-feature regime: frozen encoder features are precomputed
once, and the fusion module and heads are trained on the cache; the best
configuration is optionally fine-tuned end-to-end with LoRA adapters.
Optimization uses AdamW (learning rate $10^{-4}$ for fusion/heads, cosine
schedule, 30--50 epochs). Augmentation includes video re-compression, blur, and
color jitter, plus audio noise (MUSAN), reverberation, and codec simulation.
\todo{Finalize hyperparameters and GPU-hours after experiments.}

\section{Evaluation Metrics}
\begin{itemize}
  \item \textbf{Detection:} AUC (primary), EER, accuracy, macro-F1 --- reported
        per modality and for the four-class quadrant task.
  \item \textbf{Localization:} AP@IoU $\{0.5, 0.75, 0.95\}$ and AR@$\{10,20,50\}$,
        per modality.
  \item \textbf{Generalization:} cross-dataset AUC and LOGO AUC.
  \item \textbf{Calibration:} Expected Calibration Error and reliability diagrams.
  \item \textbf{Robustness:} AUC versus compression level (H.264 CRF) and
        audio SNR.
  \item \textbf{Fairness:} AUC/EER gaps across gender, apparent skin tone, and
        language subgroups.
\end{itemize}
Statistical rigor: three seeds (mean $\pm$ std), bootstrap 95\% confidence
intervals on AUC, DeLong tests for AUC differences, and Holm--Bonferroni
correction across ablation families.

\section{Baselines}
Unimodal video: Xception~\cite{rossler2019faceforensics}, EfficientNet-B4,
LipForensics~\cite{haliassos2021lipforensics}. Unimodal audio:
RawNet2~\cite{tak2021rawnet2}, AASIST~\cite{jung2022aasist}. Multimodal:
MDS~\cite{chugh2020mds}, Emotions Don't Lie~\cite{mittal2020emotions},
AVoiD-DF~\cite{yang2023avoiddf}, and BA-TFD+~\cite{cai2022lavdf} for
localization. All baselines are evaluated under identical splits and
preprocessing.
